\documentclass[../../main.tex]{subfiles}
\begin{document}

\begin{center}
  {\fontsize{48pt}{58pt}\selectfont \textbf{1975}\normalsize\textbf{年度}}
\end{center}

\begin{table}[htb]
  \centering
  \caption{出題テーマと難易度}
  \begin{tabular}{|c|p{9cm}|c|}
    \hline
    \textbf{問題番号} & \textbf{テーマ} & \textbf{難易度} \\
    \hline
    第1問 & 三角形の面積を二等分する最短線分 & \\
    \hline
    第2問 & 多項式の恒等式と整数の組 & \\
    \hline
    第3問 & 直線に接する放物線と囲まれた面積 & \\
    \hline
    第4問 & 平方根の漸化式と三角関数 & \\
    \hline
    第5問 & 球の共通部分の体積と無限級数 & \\
    \hline
    第6問 & 玉を入れ替える操作と期待値 & \\
    \hline
  \end{tabular}
\end{table}

\begin{table}[htb]
  \centering
  \caption{解答}
  \renewcommand{\arraystretch}{1.4}
  \begin{tabular}{|c|c|p{9cm}|}
    \hline
    \textbf{問題番号} & \textbf{小問} & \textbf{答え} \\
    \hline
    第1問 & -- & $P$，$Q$は辺$BC$，$CA$上で，ともに$C$からの距離が$24$ \\
    \hline
    第2問 & -- & $(k,l,m,n)=(1,1,0,1)$ \\
    \hline
    第3問 & -- & $(a,b)=(1,3)$のとき，最大面積は$\dfrac{9}{2}$ \\
    \hline
    第4問 & -- & $\theta_n=\dfrac{\pi}{2}-\dfrac{\pi}{4}\left(\dfrac{1}{2}\right)^{n-1}$，極限値は$\dfrac{\pi}{2}$ \\
    \hline
    第5問 & (1) & $v_n=\dfrac{104\pi a^3}{3\cdot2^{3n+9}}$ \\
    \cline{2-3}
    & (2) & $V=\dfrac{13\pi a^3}{1344}$ \\
    \hline
    第6問 & (1) & 与えられた漸化式が成立（証明問題） \\
    \cline{2-3}
    & (2) & $E_n=3-2\left(\dfrac{5}{6}\right)^n$ \\
    \cline{2-3}
    & (3) & 極限値は$3$ \\
    \hline
  \end{tabular}
\end{table}

\paragraph{試験形態}
6問．試験時間は要確認．

\end{document}
